\chapter{Introduction}

The objective of this document is to introduce the Final Project for the \textit{Electronics for Embedded Systems} course at Politecnico di Torino, developed by Andrea Viola. The main goal is to describe the context and motivation behind the project and, from these, define a first set of specifications to be met by its final working implementation.

In particular, the project has been selected to provide the opportunity to explore an application in the automotive field, which is the area of electronics I am most interested in.

The proposed project consists of the design and implementation of a \textbf{Steer-by-Wire automotive system demonstrator}. A Steer-by-Wire system is, by nature, a strongly safety-constrained embedded system which, in a real automotive application, must comply with strict requirements, development processes and safety regulations.

Within the scope and time constraints of an academic project, a deliberately simplified version of the electronic architecture will be designed and implemented. The objective is not to reproduce a production-ready automotive system, but rather to demonstrate a good understanding of its main requirements, architectural implications and safety-oriented design practices. This approach also allows most of the fundamental and advanced embedded electronics concepts introduced throughout the course to be applied within a single, well-defined engineering problem, rather than being treated as isolated functional blocks.

The architecture described in the following chapters should therefore be considered as a first proposal for the project specifications. Its structure and individual subsystems may be adapted or extended according to the feedback received, in particular to better emphasize the main topics and learning objectives required for the final project.